% Generated draft: requires mathematical and source review.
\documentclass[11pt]{article}
\usepackage[margin=25mm]{geometry}
\usepackage{fontspec}
\setmainfont{Latin Modern Roman}
\usepackage{amsmath,amssymb,mathtools,unicode-math}
\setmathfont{Latin Modern Math}
\usepackage{enumitem,needspace,xurl,hyperref}
\setlength{\parindent}{0pt}
\setlength{\parskip}{5pt}
\setlength{\emergencystretch}{3em}
\newcommand{\N}{\mathbb{N}}
\newcommand{\Z}{\mathbb{Z}}
\newcommand{\Q}{\mathbb{Q}}
\newcommand{\R}{\mathbb{R}}
\newcommand{\C}{\mathbb{C}}
\newcommand{\F}{\mathbb{F}}
\newcommand{\A}{\mathbb{A}}
\newcommand{\PP}{\mathbb{P}}
\newcommand{\E}{\mathbb{E}}
\newcommand{\eps}{\varepsilon}
\newcommand{\dd}{\,\mathrm{d}}
\begin{document}
\section*{UnsolvedMath ID: 800140}
\textbf{Source problem code:} catalogue entry
\subsection{Definitions and mathematical statement}
The supplied record contains two separately attributed passages. The first gives a qualitative problem; the second ends with an asserted asymptotic and supplies no further question.

Alexandre Bailleul (Université Paris--Saclay) asks about Chebyshev's bias over function fields and the arguments of zeros of the associated \(L\)-functions. The source does not specify the function fields, the family of \(L\)-functions, the selected zeros, or conventions for their arguments. For a finite selected family of real arguments \(\theta_1,\ldots,\theta_r\), linear independence over \(\Q\) means
\[
\sum_{j=1}^{r}q_j\theta_j=0,\quad q_j\in\Q
\quad\Longrightarrow\quad q_1=\cdots=q_r=0.
\]
Linear independence together with \(\pi\) means
\[
q_0\pi+\sum_{j=1}^{r}q_j\theta_j=0,\quad q_0,\ldots,q_r\in\Q
\quad\Longrightarrow\quad q_0=\cdots=q_r=0.
\]
Under the first independence assumption, failure of the second is equivalent to the existence of \(c_1,\ldots,c_r\in\Q\) such that
\[
\pi=\sum_{j=1}^{r}c_j\theta_j.
\]
An individual argument being a rational multiple of \(\pi\) is one possible instance of failure, but does not describe every possible relation involving several arguments.

The requested assertion is that failure of linear independence together with \(\pi\) always has a geometric explanation. The expression ``geometric explanation'' is not defined in the supplied source, so this request does not yet specify a precise theorem. The source names supersingularity of elliptic curves and Jacobians that are not absolutely simple as geometric conditions usually responsible for known instances. It also reports that no rational linear relations among the relevant zero arguments themselves are known. These are supplied source qualifications, not independently verified assertions about the literature. In particular, the unspecified choice of arguments must be resolved before that independence claim can be interpreted precisely.

The second passage, attributed to Stanley Yao Xiao (University of Northern British Columbia), considers a primitive positive definite integral binary quadratic form
\[
f(u,v)=Au^2+Buv+Cv^2,\qquad A,B,C\in\Z,\qquad \gcd(A,B,C)=1,
\]
where positive definiteness means \(f(u,v)>0\) for every \((u,v)\in\R^2\setminus\{(0,0)\}\). Its discriminant is
\[
D=\operatorname{Disc}(f)=B^2-4AC<0.
\]
The additional source hypothesis is
\[
f(T,1)\not\equiv T(T+1)\pmod 2,
\]
with polynomial congruence understood in \(\F_2[T]\). For real \(X\geq2\), define
\[
\pi_f(X)=\#\bigl\{p\leq X:\ p\text{ is a positive prime and }p=f(a,b^2)\text{ for some }a,b\in\Z\bigr\}.
\]
Each prime is counted once, irrespective of its number of representations. The supplied passage attributes to the main theorem of the unidentified reference [X] the asymptotic
\[
\pi_f(X)\sim C(D)\frac{X^{3/4}}{\log X}\qquad(X\to\infty),
\]
where \(C(D)>0\) is stated to depend only on the discriminant. This is recorded as an attributed result, not as a new problem or an independently verified theorem. No subsequent question from Xiao is present in the supplied record.
\subsection{Short English statement}
Bailleul asks whether every failure of independence between function-field zero arguments and \(\pi\) has a geometric explanation. The relevant family of zeros and the meaning of that explanation remain unspecified. A separately attributed Xiao passage states an asymptotic for primes represented by \(f(a,b^2)\), but contains no explicit question.
\subsection{Sources}
\begin{itemize}
\item Local UnsolvedMath registry, permanent ID \(800140\), exact source problem code \texttt{catalogue entry}, title ``Comparative Prime Number Theory --- Problem 9.'' The supplied registry labels the record open; that label is not independent verification of either passage's mathematical status.
\item Hamieh, Kadiri, Martin, and Ng, \emph{Comparative Prime Number Theory Problem List} (2024), source item Problem (\#9), \url{https://arxiv.org/abs/2407.03530}. The supplied provenance records access on 2026-10-16 and extraction by \texttt{arxiv\_latex}; these are supplied provenance details, not a claim that the source was visited or independently checked for this draft.
\item The function-field question is attributed in the supplied passage to Alexandre Bailleul, Université Paris--Saclay. The source refers to work of Cha using the label [Cha], but supplies no full bibliographic entry for that reference.
\item The passage on generalized Friedlander--Iwaniec primes is attributed to Stanley Yao Xiao, University of Northern British Columbia. Its asserted asymptotic is attributed to the main theorem of Xiao's paper [X]; the supplied record gives no title, bibliographic entry, or link for [X].
\end{itemize}
\end{document}
